\documentclass[11pt,a4paper]{article}
\usepackage[utf8]{inputenc}
\usepackage{amsmath,amssymb,amsthm}
\usepackage{mathtools}
\usepackage{geometry}
\usepackage{hyperref}
\usepackage{enumitem}
\usepackage{xcolor}

\geometry{margin=2.5cm}

\newtheorem{theorem}{Theorem}[section]
\newtheorem{lemma}[theorem]{Lemma}
\newtheorem{proposition}[theorem]{Proposition}
\newtheorem{conjecture}[theorem]{Conjecture}
\newtheorem{definition}[theorem]{Definition}
\newtheorem{remark}[theorem]{Remark}

\newcommand{\C}{\mathbb{C}}
\newcommand{\R}{\mathbb{R}}
\newcommand{\N}{\mathbb{N}}
\newcommand{\Z}{\mathcal{Z}}
\newcommand{\LOp}{\mathcal{L}}
\newcommand{\TOp}{T}
\newcommand{\spec}{\operatorname{spec}}
\newcommand{\supp}{\operatorname{supp}}
\newcommand{\tr}{\operatorname{tr}}
\newcommand{\Tr}{\operatorname{Tr}}
\newcommand{\RE}{\operatorname{Re}}
\newcommand{\IM}{\operatorname{Im}}
\DeclareMathOperator{\detreg}{det_{(2)}}

\title{\bf Advancing the Iota Walk Approach to the Riemann Hypothesis\\
Exploration of Three Avenues}
\author{}
\date{\today}

\begin{document}
\maketitle

\begin{abstract}
We investigate three proposed directions for turning the geometric greedy condition (GC) on the iota walk into a proof that any zero of the Dirichlet eta function satisfying GC must lie on the critical line $\sigma = 1/2$.  The avenues are: (1) analytic estimates linking GC to exponential sums and the alternating series; (2) analysis of the continuous integral walk derived from Riemann's representation; (3) connection to the GDST correlation operator and its Fredholm determinant.  We present what can be proved with current tools, identify the precise obstacles, and formulate a series of conjectures whose resolution would complete the proof of the Riemann Hypothesis.
\end{abstract}

\section{Introduction}
We recall the iota walk for $s = \sigma + it$:
\[
v_n(s) = (-1)^{n-1} n^{-s},\qquad 
\mathbf{S}_N(s) = \sum_{n=1}^{N} v_n(s).
\]
The limit (when it exists) is the Dirichlet eta function $\eta(s)$.  A necessary and sufficient condition for a zero of $\eta(s)$ on the critical strip to lie on the line $\sigma = 1/2$ is given by the \textbf{greedy condition} (GC):
\begin{equation}\label{GC}
\langle\mathbf{S}_{N-1}(s), v_N(s)\rangle \le 0 \qquad \forall N \ge 1,
\end{equation}
where $\langle\cdot,\cdot\rangle$ is the Euclidean dot product in $\R^2$.  It is known that every computed zero on the critical line satisfies GC; conversely, if a zero satisfies GC then it must lie on $\sigma = 1/2$ (this is the statement to be proved).  The present note explores three strategies to prove the converse direction.

\section{Avenue 1: Exponential sum estimates with GC}
\label{sec:exp}

The dot product in \eqref{GC} can be written as
\[
\langle\mathbf{S}_{N-1}(s), v_N(s)\rangle
= \sum_{n=1}^{N-1} (-1)^{n-1} n^{-s} \cdot (-1)^{N-1} N^{-s} \;\; \text{(real part)}.
\]
Explicitly,
\[
X_{N-1} x_N + Y_{N-1} y_N
= \RE \Bigl( \overline{\mathbf{S}_{N-1}(s)}\, v_N(s) \Bigr)
= \RE \Bigl( (-1)^{N-1} N^{-s} \sum_{n=1}^{N-1} (-1)^{n-1} n^{-\bar s} \Bigr).
\]
The condition $\le 0$ is therefore a sign restriction on the product of the $N$-th step with the complex conjugate of the current partial sum.

A first observation is that the squared distance evolves as
\[
|\mathbf{S}_N(s)|^2 = |\mathbf{S}_{N-1}(s)|^2 + 2\langle\mathbf{S}_{N-1}(s), v_N(s)\rangle + |v_N(s)|^2.
\]
Hence GC implies $|\mathbf{S}_N|^2 \le |\mathbf{S}_{N-1}|^2 + N^{-2\sigma}$.  If additionally $\lim_{N\to\infty}\mathbf{S}_N(s)=0$ (i.e.\ $s$ is a zero of $\eta$), we can telescope the inequality:
\[
|\mathbf{S}_N|^2 \le \sum_{k=N+1}^{\infty} k^{-2\sigma} \qquad (\text{provided the sum converges}),
\]
which requires $\sigma > 1/2$.  For $\sigma \le 1/2$, the series $\sum k^{-2\sigma}$ diverges, so the inequality is only useful in a very weak sense.

For $\sigma > 1/2$, the Dirichlet series converges absolutely, and one can estimate the tail directly.  Classical results on alternating zeta series give
\[
\eta(s) - \mathbf{S}_N(s) = \sum_{k=N+1}^{\infty} (-1)^{k-1} k^{-s}.
\]
The size of this remainder is $O(N^{1/2-\sigma}\log N)$ (using the bound for exponential sums).  If $\sigma > 1/2$, the remainder tends to zero, and one can show that $\eta(s) \neq 0$ except possibly at finitely many points, because $\eta(s) = (1-2^{1-s})\zeta(s)$ and $\zeta(s)$ has no zeros for $\sigma > 1$ and none on the line $\sigma = 1$ except possibly at $s = 1$ (which is not a zero).  Thus $\eta(s) \neq 0$ for all $\sigma > 1$.  For $1/2 < \sigma \le 1$, the classical zero‑free region of $\zeta(s)$ (the Littlewood–Selberg region) guarantees that $\zeta(s) \neq 0$ for $\sigma > 1 - c/\log|t|$, so $\eta(s) \neq 0$ as well, except perhaps on the boundary $\sigma = 1/2$.  Hence, if a zero of $\eta(s)$ exists with $\sigma > 1/2$, it must lie on the line $\sigma = 1/2$ by current zero‑free region results; no extra condition is needed!

The real difficulty is the region $0 < \sigma \le 1/2$.  Here the series converges conditionally, and the limit can be zero.  GC must eliminate the possibility $\sigma < 1/2$.  How can GC accomplish this?  We propose the following approach:

\begin{definition}[Approximate Dirichlet polynomials]
For $N \ge 1$, define the complex numbers
\[
\Phi_N(s) = \frac{1}{N} \sum_{n=1}^{N} n^{1/2} \cdot n^{-s} = \frac{1}{N} \sum_{n=1}^{N} n^{1/2 - s}.
\]
These are essentially averages of the step directions scaled by the square‑root of the index.
\end{definition}

The series $\sum n^{1/2 - s}$ is related to the derivative of the partial sums of the Riemann zeta function along the critical strip.  Using the Euler–Maclaurin summation, one can express $\sum_{n=1}^{N} n^{-s}$ as an integral plus a polynomial in $N$.  The alternating series can be treated similarly after combining pairs.

A promising direction is to show that GC forces the partial sums $\mathbf{S}_N(s)$ to have a well‑defined argument that rotates slowly compared to the alternating factor $(-1)^{N-1}$.  Specifically, if $\sigma < 1/2$, the terms $n^{-s}$ grow in magnitude, so the walk is dominated by the latest terms.  One might prove that for large $N$, the vector $v_N(s)$ points roughly in the direction of $\mathbf{S}_{N-1}(s)$ with high probability (in the sense of Diophantine approximation of $\ln n$ modulo $2\pi$), causing a positive dot product infinitely often.  This would violate GC.  Conversely, if $\sigma = 1/2$, the terms have constant magnitude, and the alternating sign may be arranged so that the dot product stays non‑positive.

Making this precise requires a deep study of the exponential sum $\sum_{n=1}^{N} (-1)^{n-1} n^{-it}$.  Known estimates (e.g., Vinogradov's method) give bounds like $O(\sqrt{N})$ for this sum, but the sign of the dot product is a more subtle correlation.  We may need to examine the derivative of the argument of $\mathbf{S}_{N-1}(s)$ with respect to $N$.

\begin{conjecture}[Phase‑locking]\label{conj:phase}
If $s = \sigma + it$ with $0 < \sigma < 1/2$, then for any $\varepsilon > 0$, there exist infinitely many $N$ such that
\[
|\langle\mathbf{S}_{N-1}(s), v_N(s)\rangle| \ge (1-\varepsilon) |\mathbf{S}_{N-1}(s)|\,|v_N(s)|,
\]
i.e., the step is nearly parallel to the current position.  Moreover, the dot product is positive infinitely often.
\end{conjecture}

If this conjecture holds, then GC fails for all $s$ with $\sigma < 1/2$, proving the desired implication.  The converse statement (that on $\sigma = 1/2$ the zeros satisfy GC) is supported by numerical evidence but not yet proved.

\section{Avenue 2: Analysis of the continuous walk}
\label{sec:continuous}

Riemann's integral representation for $\zeta(s)$ gives, for $\RE s = \beta$,
\[
\zeta(s) = \frac{1}{s-1} + \frac12 + 2 \int_{0}^{\infty} \frac{\sin(s \arctan u)}{(1+u^2)^{s/2} (e^{2\pi u}-1)} \,du .
\]
Setting $u = \tan\theta$, we obtain for $\beta = 1/2$ the simple form with $\sin(\theta/2)\cosh(t\theta)$ and $\cos(\theta/2)\sinh(t\theta)$.  For general $\beta$, the integrand involves $\sin(\beta\theta)\cosh(t\theta)$ and $\cos(\beta\theta)\sinh(t\theta)$, with an additional phase factor.  The continuous walk is defined by
\[
S_\beta(\alpha) = 2\int_{0}^{\alpha} \frac{\sin((\beta+it)\theta)}{(\cos\theta)^{2}(e^{2\pi\tan\theta}-1)} \,d\theta - \frac32,
\qquad 0 \le \alpha \le \frac{\pi}{2}.
\]

The continuous greedy condition is $\frac{d}{d\alpha} |S_\beta(\alpha)|^2 \le 0$ for all $\alpha$.  This is equivalent to $\RE( \overline{S_\beta(\alpha)} S_\beta'(\alpha) ) \le 0$.

Differentiating under the integral sign,
\[
S_\beta'(\alpha) = \frac{2\sin((\beta+it)\alpha)}{(\cos\alpha)^{2}(e^{2\pi\tan\alpha}-1)} .
\]
Thus the condition becomes
\[
\RE \Bigl( \overline{S_\beta(\alpha)} \cdot \frac{2\sin((\beta+it)\alpha)}{(\cos\alpha)^{2}(e^{2\pi\tan\alpha}-1)} \Bigr) \le 0 \qquad \forall \alpha \in [0,\pi/2].
\]

Using the fact that the denominator is real and positive, this reduces to
\[
\RE \bigl( \overline{S_\beta(\alpha)} \sin((\beta+it)\alpha) \bigr) \le 0.
\]

Now, $\sin((\beta+it)\alpha) = \sin(\beta\alpha)\cosh(t\alpha) + i \cos(\beta\alpha)\sinh(t\alpha)$.  Writing $S_\beta(\alpha) = A(\alpha) + i B(\alpha)$, the inequality is
\[
A(\alpha) \sin(\beta\alpha)\cosh(t\alpha) + B(\alpha) \cos(\beta\alpha)\sinh(t\alpha) \le 0.
\tag{2}
\]

At $\alpha = 0$, $S_\beta(0) = -3/2$, $A(0)=-3/2$, $B(0)=0$, $\sin(0)=0$, $\cos(0)=1$, so the left‑hand side is $(-3/2) \cdot 0 \cdot \cosh(0) + 0 \cdot 1 \cdot \sinh(0) = 0$, so the condition holds at the start.

For small $\alpha$, expand $S_\beta(\alpha)$ in powers of $\alpha$ using the Taylor series of the integrand.  One finds
\[
S_\beta(\alpha) = -\frac32 + 2 \int_{0}^{\alpha} \frac{ (\beta+it)\theta + O(\theta^2) }{ (1+O(\theta^2)) ( 2\pi\theta + O(\theta^2) ) }\,d\theta
= -\frac32 + \frac{\beta+it}{\pi} \int_{0}^{\alpha} \frac{d\theta}{\theta + O(\theta^2)} + \cdots
\]
The exact leading behaviour is not essential; the key is that the derivative $S_\beta'(0)$ is finite and given by a limit.  The inequality (2) at leading order in $\alpha$ yields a condition relating $\beta$ and $t$.  In particular, at $\alpha = 0$ the second derivative of $|S_\beta|^2$ is
\[
\frac{d^2}{d\alpha^2} |S_\beta(\alpha)|^2 \Big|_{\alpha=0} = 2|S_\beta'(0)|^2 + 2\RE( \overline{S_\beta(0)} S_\beta''(0) ),
\]
so if this is negative, the condition holds.

A full analytic treatment of the global inequality (2) is challenging, but we may attempt to prove that it can hold only when $\beta = 1/2$ by using the functional equation.  Observe that the integral representation is valid only for $\RE s = \beta$; the continuation to the whole plane depends on the functional equation, which might impose symmetry constraints on $S_\beta$.  In particular, the functional equation relates $S_{1-\beta}$ to $S_\beta$; if the walk for $\beta$ is monotone decreasing and ends at zero, then the walk for $1-\beta$ would also end at zero, and the monotonicity would have to reverse.  This might force $\beta = 1/2$.

\begin{conjecture}[Continuous monotonicity forces symmetry]\label{conj:cont}
If $S_\beta(\pi/2) = 0$ and the continuous greedy condition holds, then necessarily $\beta = 1/2$.
\end{conjecture}

This avenue remains largely speculative; a rigorous proof would likely require a very detailed analysis of the integral, possibly using the Fourier expansion of the kernel $1/(e^{2\pi\tan\theta}-1)$ and the bound of exponential sums.

\section{Avenue 3: Connection to the GDST operator}
\label{sec:GDST}

The GDST programme constructed a correlation operator $T_\sigma$ on $L^2[0,\pi]$ with a Fredholm determinant identity
\[
\det\nolimits_{(2)}(I - \mathcal{L}_s) = \frac{\zeta(s)}{\zeta(2s)} e^{P(s)},
\]
where $\mathcal{L}_s$ is a transfer operator unitarily related to $T_\sigma$.  The eigenvalues $\lambda_i(\sigma)$ of $T_\sigma$ are real.  The fatal mistake of that programme was the claimed identity
\[
\Tr\bigl[ (z - T_\sigma)^{-1} \bigr] = \sum_{\rho} \frac{1}{z - (\rho-\sigma)} + \text{entire},
\]
which would force the zeros to be real.  This identity does not hold.

However, the iota walk and its greedy condition may provide an independent link to the operator $T_\sigma$.  Consider the sequence of partial sums $\mathbf{S}_N(s)$ for $s = \sigma + it$.  The vector $\mathbf{S}_N(s)$ is essentially a linear combination of the complex numbers $n^{-s}$ with coefficients $\pm 1$.  The correlation operator $T_\sigma$ was built from the binary digits $\delta_n(\theta)$, which encode the greedy expansion of $\theta$.  If we set $\theta$ such that $E(\theta,s)$ (the GDST sub‑sum) approximates the eta function, perhaps the greedy condition on the iota walk translates into a positivity property of the operator $T_\sigma$ restricted to the subspace corresponding to the selected indices.

Specifically, the iota walk uses the alternating series, which can be expressed as the sum over all indices of the eta function.  In the GDST framework, the full Dirichlet series is split into two sub‑sums based on a parameter $\theta$:
\[
\zeta(s) - 1 = E(\theta,s) + \Omega(\theta,s),
\]
where $E(\theta,s)$ runs over indices selected by the greedy algorithm at angle $\theta$.  At a zero of $\eta(s)$ (hence of $\zeta(s)$), both $E$ and $\Omega$ must cancel.  The greedy condition in the iota walk is a statement about the partial sums of the alternating series; perhaps it corresponds to the requirement that the selected sub‑sum $E$ never dominates the omitted sub‑sum $\Omega$ in a certain sense.

If one could prove that the greedy condition implies that the spectral measure $\mu_{1/2}$ of $T_{1/2}$ has a particular simple form (e.g., it is supported on $[0,1]$ with a specific density), and then use the determinant identity to deduce that the zeros of $\zeta$ are on the critical line, one would have a proof.  But the connection between the iota walk and the operator $T_\sigma$ remains obscure.

\begin{conjecture}[GC implies a spectral gap for $T_{1/2}$]\label{conj:gap}
If $s = 1/2 + it$ is a zero of $\zeta(s)$ that satisfies the iota greedy condition, then $1$ is not an eigenvalue of $T_{1/2}$ (or $T_{1/2}$ has a spectral gap near $1$), which would force the associated Fredholm determinant to be non‑zero, contradicting the determinant identity unless $s$ is a zero on the line.  This is too vague to be useful.
\end{conjecture}

Given the failure of the original GDST spectral identification, extreme caution is needed.  Any attempt to link the iota walk to the operator $T_\sigma$ must avoid the step of equating poles of a Stieltjes transform with shifted zeros.  The iota walk provides a discrete sequence of vectors; perhaps one can construct a Jacobi matrix from the squared distances of this walk, as suggested earlier, and then show that the presence of a zero eigenvalue (corresponding to convergence to zero) implies $\sigma = 1/2$ via the asymptotics of the partial sums.  This reduces to the exponential sum estimates of Avenue~1.

\section{Conclusion}
All three avenues ultimately point to the need for deep analytic estimates of exponential sums associated with the Dirichlet eta series, combined with an intricate sign restriction (the greedy condition).  The continuous analogue may offer a more systematic approach, but it remains largely unexplored.  The connection to the GDST operator, while tempting, is fraught with the same pitfalls that derailed the original programme.

We summarise the current status:

\begin{center}
\begin{tabular}{|p{3.5cm}|p{3.5cm}|p{3.5cm}|p{2.5cm}|}
\hline
\textbf{Avenue} & \textbf{Key difficulty} & \textbf{Possible progress} & \textbf{Outlook} \\
\hline
Exponential sums + GC & Sign of dot product over large $N$ & Phase‑locking conjecture for $\sigma<1/2$ & Very hard \\
\hline
Continuous greedy condition & Global differential inequality & Small‑$\alpha$ expansion; symmetry from functional equation & Promising but unproven \\
\hline
GDST operator connection & Avoid fatal identification; find correct link & Jacobi matrix from iota walk; reduce to exponential sums & Highly uncertain \\
\hline
\end{tabular}
\end{center}

At present, none of these avenues yields a proof of the Riemann Hypothesis.  The research programme continues.

\begin{thebibliography}{9}
\bibitem{Titchmarsh} E.\,C.\ Titchmarsh, \textit{The Theory of the Riemann Zeta‑function}, 2nd ed., Oxford Univ.\ Press, 1986.
\bibitem{Iota} V.\ Geere, \textit{On the Iota Function and the Greedy Condition: A Reformulation, Not a Proof}, manuscript, 2026.
\bibitem{GDST} V.\ Geere et al., \textit{The Greedy Dirichlet Sub‑Sum Transform and the Riemann zeta function}, preprint, 2026.
\end{thebibliography}

\end{document}